\section{[Scratch] $\mathbb F2\mathbb Z$ -- general version}
\ulrich{%
How about the math letters in the name? I like it better.
}

% \left\{\begin{align} {#1} \end{align}\left| \begin{align} {#2} \end{align} \right.\right\}

Let $\codedim\geq 0$ be a dimension parameter, and $\codelength\geq 0$ a blocklength parameter. Let $\FF$ be a field (not necessarily a prime field), and let $\cC\subseteq \FF^{\codelength}$ be a linear code of dimension $k$ and blocklength $\codelength$ over $\FF$.  Let $\inplen\geq0$ be a public input length parameter, and let $Q\in \FF[Y_1, \ldots, Y_{\inplen}, Z_1, \ldots, Z_{\codedim}]$  be a constraint polynomial with coefficients in $\FF$ and on $\inplen+\codedim$ variables. Consider the following relation \albert{what name?}:
%
$$
\Relation_{\cC, \FF}^{Q} \defeq \left\{ \left(\begin{aligned} &\inp=\yy\in \FF^\inplen, \\&\oracle = f \in \FF^\codelength, \\ &\wit= \bff\in \FF^\codedim\end{aligned}\right) \;\middle|\;  \begin{aligned} &\Enc_{\cC}(\bff) = f\\ \wedge \  &Q(\yy, \bff) = 0  \end{aligned} \right\}
$$
When $Q$ and $\yy$ are such that $\yy=(\rr,\mu)\in \FF^{\log\codedim}\times \FF$ and $$Q(\yy,\ff)= \mle{\bff}(\rr)-\mu,$$ we denote $\Relation_{\cC, \FF}^{Q}$  by $\Relation_{\cC, \FF}^{\mathsf{mle}}$. 


Next let $\comfield$ be a finite field of characteristic $q$,  and let $\evalfield$ be either a finite field or $\QQ$. For intuition, $\comfield$ is the field over which we will commit our witness, and $\evalfield$ is the field over which we prove linear relations on the witness. We introduce the following short-hand notation:
%
$$
\psicomfield \defeq \psi_{\comfield}: \ZZ[X] \to \comfield, \qquad \psievalfield \defeq \psi_{\evalfield}: \ZZ[X] \to \evalfield.
$$
\ulrich{%
Shouldn't be these the projections from $\mathbb Z$ onto the "integers" in the field?
} \albert{Most of the time, yes. I wanted to allow for arbitrary fields, not just prime fields, in which case the lift from the evaluation field $\evalfield$ needs to be to $\ZZ[X]$. But this makes everything hard to read}
\ulrich{%
Ah, get it you want to lift from $\mathbb F_{q^k}$. 
} \albert{Yes, but later below I just assume that $\evalfield$ is prime for simplicity
}
\ulrich{%
Then how about starting simple? One can always do remarks on the fully general scope?
In the end, a matter of taste.
In any case, if we lift to $\mathbb Z[X]$, shouldn't we have the witness $\mathbb w$ below also in $\mathbb Z[X]^\ell$?
}
\albert{Yes, better to start simple. I'm worried the current description is getting too complicated}
\ulrich{We can always generalize later on. 
Maybe it turns out the way to go, but I think starting simple is always better}
\ulrich{%
Are we even sure that we want to do it in full generality, like arbitrary commit field... 
If it is another prime field... let me think
Anyways, need to go for check with Bowie to the vet, will return in about 2 hours.
}
\albert{Maybe it's better to fix the commit field to a binary field. Then have some remark or appendix (or nothing) with fully general commit and eval fields. Good luck at the vet!}
%
Let $\cC\subseteq \comfield^{n}$ be a linear code of dimension $k$ and blocklength $n$ over $\comfield$. Define \albert{can we handle more general witness subsets?}
$$
\Relation_{\cC, \comfield, \evalfield}^{\mathsf{lin}} \defeq \left\{ \left(\begin{aligned} &\inp=(\mu, \vv)\in \evalfield \times \evalfield^{\codedim}, \\&\oracle = f \in \comfield^\codelength, \\ &\wit= \bff\in \range{0}{q-1}^\codedim \subseteq \ZZ^\codedim\end{aligned}\right) \;\middle|\;  \begin{aligned} &\Enc_{\cC}(\psicomfield(\bff)) = f\\ \wedge \  & \langle \psievalfield(\bff), \vv\rangle  = \mu  \end{aligned} \right\}
$$
We provide an IOPP for $\Relation_{\cC, \comfield, \evalfield}^{\mathsf{lin}}$ next.


\paragraph{Protocol assumptions} Assume for simplicity that $\evalfield$ is a prime field. The general case (when $\evalfield$ is a non-trivial field extension or the field of rationals) is handled in \cref{?}.


Fix an element $\basegenerator\in \comfield$ with sufficiently large multiplicative order, to be determined later.  Let  $S_q=\psicomfield(\range{0}{q-1})$, and define the following bivariate polynomial with coefficients in $\comfield$:
$$
P(Y,Z) = \sum_{s\in \range{0}{q-1}} L_{ S_q, \psicomfield(s)}(Z) \cdot Y^{s} \in \comfield[Y,Z],
$$
where $L_{S_q, \psicomfield(s)}(Z)$ is the Lagrange polynomial associated to $\psicomfield(s)$ over the set $S_q$, with degree at most $q-1$. Note
%
$$
P(Y,\psicomfield(s)) = Y^{s} \quad \forall \ s\in \range{0}{q-1}.
$$
%
We assume prover and verifier have access to an IOR $\prod_{\Relation_{\cC, \comfield}^{Q}  \to \Relation_{\cC, \comfield}^{\mathsf{eval}}}$ \albert{define in prelims} that reduces 
%
$
\Relation_{\cC, \comfield}^{Q}$ to $\Relation_{\cC, \comfield}^{\mathsf{eval}}$, 
where $$Q(Y_1, \ldots, Y_{\codedim+1}, Z_1, \ldots, Z_{\codedim}) = \prod_{i=1}^{\codedim} P(Y_i, Z_i) - Y_{\codedim+1}.$$ We further assume that prover and verifier have access to a batched version of $\prod_{\Relation_{\cC, \comfield}^{Q}  \to \Relation_{\cC, \comfield}^{\mathsf{eval}}}$ (which allows to reduce several instances of ), and to an IOPP for $\Relation_{\cC, \comfield}^{\mathsf{eval}}$.


\begin{construction}[Core IOP for $\Relation_{\cC, \comfield, \evalfield}^{\mathsf{lin}}$]\label{c:core_iop}
The prover and verifier receive $\inp=(\mu, \vv)\in \evalfield \times \evalfield^{\codedim}$ and  $\oracle= f \in \comfield^n$. The prover receives $\wit= \bff\in \range{0}{q-1}^\codedim \subseteq \ZZ^\codedim$ such that $(\inp,\oracle, \wit) \in \Relation_{\cC, \comfield, \evalfield}^{\mathsf{lin}}$. 

Fix a tensor decomposition of $\vv$ as $$\vv = \vv^{(1)} \otimes \vv^{(2)} = \left(\vv_i^{(1)}\cdot \vv_j^{(2)}\right)_{i\in [\codedim_1], j\in [\codedim_2]}$$ for some $\vv^{(1)}\in \evalfield^{\ell_1}, \vv^{(2)}\in \evalfield^{\ell_2}$, with $\ell_1 \cdot \ell_2=\ell$. Note that the trivial tensor decomposition is admitted, i.e.\ we may not decompose $\vv$ and set $\vv^{(1)}=\vv$ and $\vv^{(2)} = (1)$.

Index the entries of $\bff$ as $\bff=(\bff_{ij})_{i\in [\ell_1], j\in[\ell_2]}$. Note that
%
$$
\langle \psievalfield(\bff), \vv \rangle = \sum_{j\in [\ell_2]}\left(\sum_{i\in [\ell_1]} \psievalfield(\bff_{ij}) \cdot \vv_i^{(1)}\right) \cdot \vv_j^{(2)} = \sum_{j\in [\ell_2]} \langle \psievalfield(\bff_{*j}) , \vv^{(1)} \rangle \cdot \vv_j^{(2)},
$$
where $\bff_{*j} = (\bff_{ij})_{i\in[\ell_1]}$.

\begin{enumerate}[leftmargin=*]
\item The prover sends $\mu_1', \ldots, \mu_{\ell_2}'\in \ZZ$ such that $\langle \bff_{*j}, \psievalfield^{-1}(\vv^{(1)})\rangle = \mu_j'$. 
\item The verifier checks that $\sum_{j\in [\ell_1]} \psievalfield(\mu_j')\cdot \vv_j^{(2)}  = \mu$. 
\item Prover and verifier define $\bm{\gamma}_i = \psievalfield^{-1}(\vv_i^{(1)})$ for $i \in [\ell_1]$, and set to prove the following claim:
%
$$
\prod_{i \in [\codedim]} g^{\bm{\gamma}_i\cdot \bff_{ij}}  = g^{\mu_j'} \qquad \text{for all } j\in [\ell_2].
$$
Note that this claim is equivalent to
$$
\prod_{i \in [\codedim]} P\left(g^{\bm{\gamma}_i}, \psicomfield(\bff_{ij})\right) = \prod_{i \in [\codedim]} \sum_{s\in \range{0}{q-1}} L_{ S_q, \psicomfield(s)}(\psicomfield(\bff_{ij})) \cdot g^{\bm{\gamma}_i\cdot s}= g^{\mu_j'}  \qquad \text{for all } j\in [\ell_2].
$$
To prove the above claim, prover and verifier engage in the batched version of the IOR $\prod_{\Relation_{\cC, \comfield}^{Q}  \to \Relation_{\cC, \comfield}^{\mathsf{eval}}}$ with inputs $$\left(\inp'=((g^{\bm{\gamma}_i})_{i \in [\codedim_1]}), g^{\mu_j'}), \oracle_j'=f_j, \wit_j'=\psicomfield(\bff_{*j})\right)  \qquad \text{for } j\in [\ell_2].$$
%
\item At the end of this IOR, they are left with the following claim to prove:
%
$$
((\mu', \rr'), f', \bff') \in \Relation_{\cC, \comfield}^{\mathsf{eval}}.$$ Accordingly, the prover and the verifier engage in a IOPP  to certify the validity of this last claim.
\end{enumerate}
\end{construction}

. 








\section{Scratch}


\begin{construction}
The verifier receives $\inp=(\mu, \vv)\in \evalfield \times \evalfield^{\codedim}$ and oracle access to $\oracle= f \in \comfield^n$. The prover receives $\wit= \bff\in \range{0}{q-1}^k \subseteq \ZZ^k$ such that $(\inp,\oracle, \wit) \in \Relation_{\cC, \comfield, \evalfield}^{\mathsf{lin}}$. 


\begin{enumerate}[leftmargin=*]
\item The prover sends $\mu'\in \ZZ[X]$ such that $\langle \bff, \psievalfield^{-1}(\vv)\rangle = \mu'$. 
\item The verifier checks that $\psievalfield(\mu') = \mu$. Then, the verifier samples $\alpha\in [0, 2^N-1]\subseteq \ZZ$ and sends $\alpha$ to the prover. Let $\psi_\alpha: \ZZ[X]\to \ZZ$ be the evaluation map at $\alpha$, so that $\psi_\alpha(a(X))= a(\alpha)$. 
\item The prover and verifier engage in a PIOP to prove the following equality:
%
$$
\prod_{i \in [\codedim]} g^{\psi_\alpha(\psievalfield^{-1}(\vv_{i}))\cdot \bff_i}  = g^{\psi_\alpha(\mu')}.
$$
\end{enumerate}
\end{construction}





\begin{enumerate}[leftmargin=*]
\item The prover sends $\mu'\in \ZZ$ such that $\langle \bff, \psievalfield^{-1}(\vv)\rangle = \mu'$. 
\item The verifier checks that $\psievalfield(\mu') = \mu$. 
\item Prover and verifier define $\bm{\gamma}_i = \psievalfield^{-1}(\vv_i)$ for $i \in [\codedim]$, and set to prove the following claim:
%
$$
\prod_{i \in [\codedim]} g^{\bm{\gamma}_i\cdot \bff_i}  = g^{\mu'}.
$$
Note that this claim is equivalent to
$$
\prod_{i \in [\codedim]} P\left(g^{\bm{\gamma}_i}, \psicomfield(\bff_i)\right) = \prod_{i \in [\codedim]} \sum_{s\in \range{0}{q-1}} L_{ S_q, \psicomfield(s)}(\psicomfield(\bff_i)) \cdot g^{\bm{\gamma}_i\cdot s}= g^{\mu'}.
$$
To prove the above equality, prover and verifier engage in the IOR $\prod_{\Relation_{\cC, \comfield}^{Q}  \to \Relation_{\cC, \comfield}^{\mathsf{eval}}}$ with inputs $$\left(\inp'=((g^{\bm{\gamma}_i})_{i \in [\codedim]}), g^{\mu'}), \oracle'=f, \wit'=\psicomfield(\bff)\right).$$
%
\item At the end of this IOR, they are left with the following claim to prove:
%
$$
((\mu', \rr'), f', \bff') \in \Relation_{\cC, \comfield}^{\mathsf{eval}}.$$ Accordingly, the prover and the verifier engage in a IOPP  to certify the validity of this last claim.
\end{enumerate}



\appendix
\printbibliography[heading=bibintoc,title={References}]

